%==============================================================================
%  MATH 231 -- Linear Algebra I (Queens College, Fall 2026)
%  HOMEWORK 1 -- Vectors and Complex Numbers
%
%  SCOPE.  This file currently covers exactly two lectures:
%     lectures/vector-unit/MATH231_vectors_1_definition_and_arithmetic  (S1-S9)
%     lectures/vector-unit/MATH231_complex_numbers                     (S1-S7)
%  The complex-numbers lecture is deliberately cut BEFORE S8 (de Moivre, the
%  roots of unity, the fundamental theorem of algebra) and before the quaternion
%  material -- none of that has been covered in class yet.
%
%  MORE PROBLEMS TO COME.  Parts are lettered (A, B, C, D, E) rather than
%  numbered, and problems are numbered by part (Problem A.1, B.2, ...), so
%  further parts from later lectures can be appended without renumbering
%  anything already written.  Add them as new \section blocks before Part E, or
%  extend Part E's Section E.x list, whichever fits.
%
%  DELIBERATE OMISSIONS, so that nothing here outruns the lectures:
%   - No dot product.  It arrives in vectors Part 2.  The complex lecture's S7
%     proof that multiplication by i is a rotation *does* use it, so Problem C.7
%     and D.5 are routed through the modulus instead, and the forward reference
%     is flagged honestly in the text.
%   - No unit vectors / normalisation (\hat{u}), no S^1 notation, no metrics.
%   - No argument/polar form beyond cos(theta) + i sin(theta) and Euler.
%
%  NO ANSWER KEY, by design.  Part E is the answer key: every arithmetic problem
%  in Parts A and C can be checked by the student in NumPy.  That is the whole
%  pedagogical point of the programming section, and it is stated as such.
%
%  TWO THINGS TO SET BEFORE HANDING THIS OUT:
%     \duedate  -- currently a placeholder.
%     The submission paths in the "What to submit" callout, if you want a
%     different directory layout than homework/hw1/ in the student repo.
%
%  Compile with pdflatex (standard TeX Live packages only; uses listings).
%==============================================================================

\documentclass[11pt]{article}

\usepackage[margin=1in]{geometry}
\usepackage{amsmath,amssymb,amsthm}
\usepackage[dvipsnames]{xcolor}
\usepackage{enumitem}
\usepackage{booktabs}
\usepackage{array}
\usepackage[hidelinks]{hyperref}
\usepackage{titlesec}
\usepackage{needspace}
\usepackage{listings}

%------------------------------------------------------------------ style ----
\titleformat{\section}{\large\bfseries\sffamily\color{RoyalBlue}}
                      {Part \thesection.}{0.6em}{}
\titleformat{\subsection}{\normalsize\bfseries\sffamily}{\thesubsection.}{0.5em}{}
\titlespacing*{\section}{0pt}{16pt}{7pt}

\renewcommand{\thesection}{\Alph{section}}

\theoremstyle{definition}
% Extra air above and below every problem, so consecutive questions read as
% separate blocks rather than one continuous column of text.
\makeatletter
\def\thm@space@setup{%
  \thm@preskip=20pt plus 6pt minus 4pt
  \thm@postskip=\thm@preskip}
\makeatother
\newtheorem{prob}{Problem}[section]

% code blocks -- same style as Setup Guide 2, so the two documents look alike
\lstdefinestyle{term}{
  basicstyle=\ttfamily\small,
  backgroundcolor=\color{gray!10},
  frame=single, rulecolor=\color{gray!45},
  framesep=5pt, xleftmargin=8pt, xrightmargin=8pt,
  breaklines=true, columns=fullflexible, keepspaces=true,
  showstringspaces=false, aboveskip=8pt, belowskip=8pt
}
\lstdefinestyle{py}{
  style=term, language=Python,
  keywordstyle=\color{RoyalBlue}\bfseries,
  commentstyle=\color{gray!80}\itshape,
  stringstyle=\color{ForestGreen},
  emph={np,numpy}, emphstyle=\color{BrickRed}
}
\lstset{style=py}

\newcommand{\heads}[1]{\par\smallskip\noindent
  {\small\sffamily\textcolor{BrickRed}{\textbf{Watch out.}}\ #1}\par\smallskip}
\newcommand{\hint}[1]{\par\smallskip\noindent
  {\small\sffamily\textcolor{RoyalBlue}{\textbf{Hint.}}\ #1}\par\smallskip}
\newcommand{\checkpoint}[1]{\par\smallskip\noindent
  {\small\sffamily\textcolor{ForestGreen}{\textbf{Checkpoint.}}\ #1}\par\smallskip}

% shaded call-out box.  The body is collected into a savebox first, because
% \colorbox needs its whole argument at once.
\definecolor{softbg}{HTML}{F2F6FA}
\newsavebox{\qcbox}
\newenvironment{callout}[1]
  {\par\medskip\needspace{5\baselineskip}%
   \begin{lrbox}{\qcbox}%
   \begin{minipage}{\dimexpr\textwidth-2\fboxsep-4pt}%
   \setlength{\parskip}{0.4em}%
   \textbf{\sffamily\color{RoyalBlue}#1}\par}
  {\end{minipage}\end{lrbox}%
   \begin{center}\colorbox{softbg}{\usebox{\qcbox}}\end{center}\medskip}
%--------------------------------------------------------------- notation ----
\newcommand{\R}{\mathbb{R}}
\newcommand{\C}{\mathbb{C}}
\newcommand{\vc}[1]{\mathbf{#1}}
\newcommand{\norm}[1]{\lVert #1\rVert}
\newcommand{\conj}[1]{\overline{#1}}
\newcommand{\Rez}{\operatorname{Re}}
\newcommand{\Imz}{\operatorname{Im}}

%--------------------------------------------------------- fill this in ------
\newcommand{\duedate}{to be announced on the class Discord}

\title{\vspace{-1.2cm}\textbf{\sffamily MATH 231 -- Linear Algebra I}\\[4pt]
       {\large\sffamily Homework 1 \textbullet\ Vectors and Complex Numbers}\\[2pt]
       {\normalsize\sffamily Kiely Hall 326 \textbullet\
        Instructor: Kevin Bedoya}}
\author{}
\date{}

\begin{document}
\maketitle
\vspace{-1.6cm}

\noindent\textbf{Due:} \duedate.

\medskip
\begin{center}
\renewcommand{\arraystretch}{1.25}
\begin{tabular}{@{}l l >{\raggedright\arraybackslash}p{7.4cm}@{}}
\toprule
\textbf{Lecture} & \textbf{Sections} & \textbf{Topics} \\
\midrule
Vector Unit, Part 1 & \S1--\S9 &
  magnitude and direction; the Euclidean norm; vector vs.\ scalar quantities;
  notation, $\R^{2}$, and signatures; addition and subtraction; the
  parallelogram rule and the triangle inequality; scalar multiplication;
  the eight algebraic laws; equality of vectors \\[3pt]
Complex Numbers & \S1--\S7 &
  the imaginary unit; $a+bi$, real and imaginary parts; the complex plane;
  addition, multiplication, the conjugate, division; the modulus, distance,
  and discs; the unit circle, trigonometric form, Euler's formula;
  multiplication by $i$ as a rotation \\
\bottomrule
\end{tabular}
\end{center}

\begin{callout}{What to submit}
Push to your own MATH 231 GitHub repository, in a folder named
\texttt{homework/hw1/}:

\begin{itemize}[leftmargin=1.4em,itemsep=7pt,topsep=3pt]
  \item \textbf{Parts A--D} --- your written work, as a single PDF named
        \texttt{hw1.pdf}. Handwritten and scanned or photographed is completely
        fine, as long as it is legible and the pages are in order.

        \textbf{PDF is the only accepted format for the written part.} A
        \texttt{.docx}, a \texttt{.jpg}, a \texttt{.png}, a \texttt{.heic}, or a
        link to a file hosted somewhere else will not be graded. Phone scanning
        apps export PDF directly, and every word processor can ``Save as PDF''.

        \textbf{Type Parts A--D up instead and you earn $+5$ points of extra
        credit on this assignment.} Any typesetting system counts --- \LaTeX,
        Word's equation editor, or anything else that produces readable
        mathematics --- and the extra credit is for the whole of Parts A--D, not
        for typing up some of it. The file you submit must still be a PDF.

  \item \textbf{Part E} --- two Python files, \texttt{hw1.py} (your work, one
        function per question part) and \texttt{runhw1.py} (which imports those
        functions and runs them all, labelling each answer). Do not submit a copy
        of your output: \texttt{runhw1.py} is \textbf{run as submitted}, and Part
        E is graded from what it prints. \emph{How to structure your submission},
        in Part E, specifies both files in full.

        \textbf{Python is the only language accepted for the programming.} Any
        other language earns no credit for Part E, however correct it is.
\end{itemize}

If you have not set up the repository yet, Setup Guide 1 (\S\,``Create your
folders'' and \S\,``Save your work and send it to GitHub'') is the walkthrough.
Only pushed work can be graded.
\end{callout}

\begin{callout}{Label every answer}
\textbf{Parts A--D.} Label each answer with the number of the problem it answers
--- \texttt{A.2(f)}, \texttt{C.5(d)}, \texttt{D.1(b)} --- and put your answers
\textbf{in the same order as they appear in this assignment}. A grader who
cannot tell which question a piece of work belongs to cannot award it credit,
and that is the most common way a correct answer loses points.

\textbf{You do not have to copy the question out.} This applies to handwritten
and typed submissions alike: the label alone is enough. Write \texttt{A.2(f)},
then your work. Do not reproduce the wording of the problem.

\textbf{Part E.} The labelling is built into the file structure specified under
\emph{How to structure your submission}: one function per question part in
\texttt{hw1.py}, named after that part, with a comment naming it directly above;
and one labelled line per part in \texttt{runhw1.py}. Code that runs correctly
but cannot be matched to a question is graded the same way an unlabelled answer
sheet is.
\end{callout}

\noindent\textbf{Grading.} Every problem is worth the same number of points.
Part E is graded by running \texttt{runhw1.py} and reading what it prints --- on
whether it produces correct answers, not on programming style.

\medskip
\noindent\textbf{Use Part E to check Parts A and C.} Every numerical answer in
those two parts can be verified in two lines of NumPy, and you are encouraged to
check all of them before submitting --- that is what the tool is for. Two limits.
The hand work must still be shown: a correct number with no derivation earns
partial credit at best. And no proof in Part B or Part D can be checked this way,
since code confirms a statement only for the numbers you tried, not for all of
them.

\medskip
\noindent\textbf{On the proof problems (Parts B and D).} These are short, and
none of them requires a clever idea. What they require is that you write down
what you are assuming, work in components or from the stated laws, justify each
equality by naming the fact that licenses it, and end with a sentence saying
what you have shown. A chain of equalities with no words is not a proof; neither
is a paragraph of words with no computation. Two or three sentences and three or
four lines of algebra is the expected length.

\medskip
\noindent\textbf{Collaboration} is encouraged --- work at a board together, argue
about the problems. What you submit must be written by you, in your own words
and your own code.

%==============================================================================
\section{Vectors --- computation}
%==============================================================================
\emph{Vector Unit, Part 1.} Graph paper helps for the drawing parts; neat
freehand axes are acceptable. Leave norms in exact form ($\sqrt{13}$,
$3\sqrt{5}$) unless a decimal is asked for --- a norm is generically irrational,
and $\sqrt{13}$ is a finished answer.

\begin{prob}[Reading a vector off an arrow]
Each row below describes an arrow in the plane by its tail and its head.

\begin{center}
\renewcommand{\arraystretch}{1.2}
\begin{tabular}{@{}c c c@{}}
\toprule
 & \textbf{Tail} & \textbf{Head} \\
\midrule
(a) & $(0,0)$   & $(5,12)$  \\
(b) & $(2,-1)$  & $(6,2)$   \\
(c) & $(-3,4)$  & $(1,7)$   \\
(d) & $(1,1)$   & $(-2,5)$  \\
\bottomrule
\end{tabular}
\end{center}

\begin{enumerate}[label=\textbf{(\roman*)},leftmargin=2.6em,itemsep=7pt,topsep=6pt]
  \item Write each arrow as a vector in bracket notation, using
        \emph{head minus tail}.
  \item Compute the norm of each.
  \item Exactly one pair among (a)--(d) consists of two arrows that are the
        \emph{same vector}. Which pair, and why? Draw both arrows on one set of
        axes to show it.
  \item Three of the four have the same norm. Explain in one sentence why that
        does \emph{not} make them the same vector, citing the definition of
        equality of vectors.
\end{enumerate}
\end{prob}

\begin{prob}[Componentwise arithmetic]
Let
\[
  \vc{u}=[\,3,\,-4\,],\qquad
  \vc{v}=[\,-1,\,2\,],\qquad
  \vc{w}=[\,5,\,0\,],\qquad
  \alpha=-3,\qquad \beta=\tfrac12 .
\]
Compute each of the following. Show the componentwise step, not just the answer.

\begin{enumerate}[label=\textbf{(\alph*)},leftmargin=2.6em,itemsep=7pt,topsep=6pt]
  \item $\vc{u}+\vc{v}$
  \item $\vc{u}-\vc{v}$
  \item $\vc{v}-\vc{u}$
  \item $\alpha\vc{u}$
  \item $\beta\vc{v}+\vc{u}$
  \item $2\vc{u}-3\vc{v}+\vc{w}$
  \item $\norm{\vc{u}}$, $\norm{\vc{v}}$, $\norm{\vc{w}}$
  \item $\norm{\vc{u}+\vc{v}}$
  \item $\norm{\alpha\vc{u}}$, and the value of
        $\lvert\alpha\rvert\,\norm{\vc{u}}$ --- do they agree?
\end{enumerate}
\end{prob}

\begin{prob}[The triangle inequality by hand]
For each pair, compute $\norm{\vc{u}+\vc{v}}$ and $\norm{\vc{u}}+\norm{\vc{v}}$,
and say whether the triangle inequality holds \emph{strictly} ($<$) or with
\emph{equality} ($=$).

\begin{center}
\renewcommand{\arraystretch}{1.2}
\begin{tabular}{@{}c c c@{}}
\toprule
 & $\vc{u}$ & $\vc{v}$ \\
\midrule
(a) & $[\,3,\,0\,]$   & $[\,0,\,4\,]$   \\
(b) & $[\,2,\,3\,]$   & $[\,4,\,6\,]$   \\
(c) & $[\,1,\,-2\,]$  & $[\,-2,\,4\,]$  \\
(d) & $[\,2,\,1\,]$   & $[\,0,\,0\,]$   \\
\bottomrule
\end{tabular}
\end{center}
\end{prob}

\begin{prob}[Reading the scalar]
Let $\vc{v}=[\,2,\,-3\,]$.

\begin{enumerate}[label=\textbf{(\alph*)},leftmargin=2.6em,itemsep=7pt,topsep=6pt]
  \item Write the vector on the same line through the origin as $\vc{v}$ that
        points the \emph{same} way and is \emph{four times} as long.
  \item Write the vector that points the \emph{opposite} way and is \emph{half}
        as long.
  \item Write the mirror image of $\vc{v}$ through the origin. What is its norm,
        and why could you have said so without computing?
  \item Is $[\,3,\,-4.5\,]$ a scalar multiple of $\vc{v}$? Is
        $[\,4,\,-6.1\,]$? Justify each answer.
  \item What does $0\vc{v}$ equal, and what happens to the direction?
\end{enumerate}
\end{prob}

\begin{prob}[Using the eight laws]
Let $\vc{u},\vc{v}\in\R^{2}$. Simplify each expression to the form
$\alpha\vc{u}+\beta\vc{v}$ for real numbers $\alpha,\beta$. At every step, name
the law (\textbf{A1}--\textbf{A4}, \textbf{S1}--\textbf{S4}) that licenses it.

\begin{enumerate}[label=\textbf{(\alph*)},leftmargin=2.6em,itemsep=7pt,topsep=6pt]
  \item $3(\vc{u}+2\vc{v}) - 2(\vc{u}-\vc{v})$
  \item $-(\vc{u}-\vc{v}) + 4\bigl(\vc{v}+\tfrac12\vc{u}\bigr)$
  \item $5(\vc{u}+\vc{v}) - 5\vc{u} - 5\vc{v}$
\end{enumerate}

\noindent Then substitute $\vc{u}=[\,1,\,-2\,]$ and $\vc{v}=[\,3,\,1\,]$ into
both the original expression and your simplified form for each of (a), (b), (c),
and confirm they agree.

\heads{None of the eight laws is a statement about \emph{subtraction}, so you
cannot cite one to justify a step taken on a subtraction sign directly. Rewrite
every subtraction as $\vc{u}+(-\vc{v})$ first.}
\end{prob}

\begin{prob}[Signatures]
Give the signature $f:A\to B$ of each operation below, using the placeholder-dot
convention, and say whether it is \emph{unary} or \emph{binary}.

\begin{enumerate}[label=\textbf{(\alph*)},leftmargin=2.6em,itemsep=7pt,topsep=6pt]
  \item the Euclidean norm on $\R^{2}$;
  \item addition of two vectors in $\R^{2}$;
  \item scalar multiplication;
  \item the operation sending a vector to its mirror image through the origin;
  \item the operation sending a pair of vectors $\vc{u},\vc{v}$ to the distance
        between their heads, that is, to $\norm{\vc{u}-\vc{v}}$.
\end{enumerate}

\noindent Then: two of the five above have the \emph{same} signature but are
different operations. Which two, and what does that tell you about how much
information a signature carries?
\end{prob}

\begin{prob}[Vector or scalar?]
Classify each quantity as a \textbf{vector quantity} or a \textbf{scalar
quantity}, and for each vector quantity say in a few words what its direction
means physically.

\begin{center}
\begin{minipage}{0.8\textwidth}
\begin{itemize}[leftmargin=1.4em,itemsep=1pt,topsep=2pt]
  \item[] mass \quad$\cdot$\quad velocity \quad$\cdot$\quad speed \quad$\cdot$\quad
          force \quad$\cdot$\quad temperature \quad$\cdot$\quad temperature
          gradient \quad$\cdot$\quad work \quad$\cdot$\quad momentum
          \quad$\cdot$\quad electric charge \quad$\cdot$\quad displacement
          \quad$\cdot$\quad distance travelled \quad$\cdot$\quad acceleration
\end{itemize}
\end{minipage}
\end{center}

\begin{enumerate}[label=\textbf{(\alph*)},leftmargin=2.6em,itemsep=7pt,topsep=6pt]
  \item A car drives once around a circular track at a constant $40$ mph. Is its
        speed constant? Is its velocity? Is it accelerating?
  \item The same car finishes where it started. What was its total
        \emph{displacement}, and what was its total \emph{distance travelled}?
        Explain why the two answers differ.
\end{enumerate}
\end{prob}

%==============================================================================
\section{Vectors --- proof}
%==============================================================================
\emph{Vector Unit, Part 1.} Every one of these is proved the same way the
lecture proves \textbf{A1}, \textbf{S1} and \textbf{A4}: write the vectors in
components, use the corresponding fact about \emph{real} numbers one coordinate
at a time, and reassemble --- except for Problem B.4, which deliberately forbids
components.

\begin{prob}[Reversing a difference]
Let $\vc{u},\vc{v}\in\R^{2}$.

\begin{enumerate}[label=\textbf{(\alph*)},leftmargin=2.6em,itemsep=7pt,topsep=6pt]
  \item Prove that $\vc{u}-\vc{v} = -(\vc{v}-\vc{u})$.
  \item Prove that subtraction is not associative by showing that
        \[
          (\vc{u}-\vc{v})-\vc{w}
          \;=\;
          \bigl[\vc{u}-(\vc{v}-\vc{w})\bigr] - 2\vc{w}
        \]
        for all $\vc{u},\vc{v},\vc{w}\in\R^{2}$, and give explicit
        $\vc{u},\vc{v},\vc{w}$ for which the two bracketings genuinely differ.
        Then state the one condition on $\vc{w}$ under which they agree, and
        prove it.
\end{enumerate}
\end{prob}

\begin{prob}[Two of the eight laws]
Let $\vc{u}\in\R^{2}$ and $\alpha,\beta\in\R$. Prove, in components:

\begin{enumerate}[label=\textbf{(\alph*)},leftmargin=2.6em,itemsep=7pt,topsep=6pt]
  \item \textbf{S2}: $(\alpha+\beta)\vc{u} = \alpha\vc{u}+\beta\vc{u}$.
  \item \textbf{S3}: $\alpha(\beta\vc{u}) = (\alpha\beta)\vc{u}$.
\end{enumerate}

\noindent For each, name the property of the \emph{real numbers} that your proof
ultimately rests on. Then explain, in one sentence, why \textbf{S1} and
\textbf{S2} are genuinely different statements rather than the same
distributive law written twice.
\end{prob}

\begin{prob}[Scaling multiplies the length by $\lvert\alpha\rvert$]
Prove that for every $\alpha\in\R$ and every $\vc{v}\in\R^{2}$,
\[
  \norm{\alpha\vc{v}} = \lvert\alpha\rvert\,\norm{\vc{v}} .
\]

\heads{At one point you will meet $\sqrt{\alpha^{2}}$. It equals
$\lvert\alpha\rvert$, \emph{not} $\alpha$. Say why in your write-up, and check
your claim against the case $\alpha=-3$, $\vc{v}=[\,1,\,0\,]$ --- if you had
written $\alpha$ there, what would the norm of $[-3,0]$ have come out to be, and
why is that impossible?}
\end{prob}

\begin{prob}[Working from the laws alone, without components]
This problem is about what the eight laws can do \emph{by themselves}. For this
problem only, you may not write a vector in components: you may use only
\textbf{A1}--\textbf{A4}, \textbf{S1}--\textbf{S4}, the uniqueness of the
additive inverse, and ordinary arithmetic on the real scalars.

Let $\vc{v}\in\R^{2}$.

\begin{enumerate}[label=\textbf{(\alph*)},leftmargin=2.6em,itemsep=7pt,topsep=6pt]
  \item Prove that $0\vc{v}=\vc{0}$. That is: the \emph{scalar} zero times any
        vector is the \emph{zero vector}.
        \hint{Start from $0\vc{v} = (0+0)\vc{v}$ and apply \textbf{S2}. You now
        have $\vc{x}=\vc{x}+\vc{x}$ for $\vc{x}=0\vc{v}$; finish by adding
        $-\vc{x}$ to both sides and using \textbf{A4}, \textbf{A2}, \textbf{A3}.}
  \item Using (a), \textbf{S2} and \textbf{S4}, prove that
        $\vc{v} + (-1)\vc{v} = \vc{0}$.
  \item Conclude that $(-1)\vc{v}=-\vc{v}$, where $-\vc{v}$ denotes the additive
        inverse guaranteed by \textbf{A4}. Say explicitly where you are using the
        fact that the additive inverse is \emph{unique} --- without it, (b) would
        only show that $(-1)\vc{v}$ is \emph{an} inverse, not \emph{the} inverse.
\end{enumerate}

\noindent Finally, in one or two sentences: why is it worth proving
$(-1)\vc{v}=-\vc{v}$ this way, when a single line in components would also do it?
\end{prob}

\begin{prob}[When the triangle inequality is an equality]
Prove: if $\vc{v}=\alpha\vc{u}$ for some scalar $\alpha\ge 0$, then
\[
  \norm{\vc{u}+\vc{v}} = \norm{\vc{u}}+\norm{\vc{v}} .
\]

\hint{Write $\vc{u}+\vc{v}=\vc{u}+\alpha\vc{u}$ and use \textbf{S2} and
\textbf{S4} to collect it into a single scalar multiple of $\vc{u}$. Then apply
Problem B.3. You will need $\lvert 1+\alpha\rvert = 1+\alpha$ --- say why the
hypothesis $\alpha\ge0$ is what makes that true.}

\medskip
\noindent Then: where exactly does your proof break down if $\alpha<0$? Give a
concrete $\vc{u}$ and $\alpha<0$ for which the two sides are unequal, and
identify which step fails.
\end{prob}

%==============================================================================
\section{Complex numbers --- computation}
%==============================================================================
\emph{Complex Numbers.} Leave answers in the form $a+bi$ with $a$ and $b$ real,
in lowest terms. A quotient is not finished until it is written that way ---
that is the whole point of the division rule.

\begin{prob}[Parts, and the plane]
For each of
\[
  z_1 = 4-3i,\qquad
  z_2 = -2+5i,\qquad
  z_3 = 6,\qquad
  z_4 = -i,\qquad
  z_5 = \tfrac12 + \tfrac32 i ,
\]
write down $\Rez(z)$ and $\Imz(z)$, and plot all five on one copy of the complex
plane, labelled.

\medskip
\noindent Then: which of the five lie on the real axis? Which on the imaginary
axis? And which of them are ordinary real numbers?

\heads{$\Imz(4-3i)$ is the real number $-3$, not $-3i$ and not $3$. Go back over
the definition of the imaginary part if that surprised you.}
\end{prob}

\begin{prob}[The four operations]
Let $z = 3-2i$ and $w=-1+4i$. Compute, showing your work:

\begin{enumerate}[label=\textbf{(\alph*)},leftmargin=2.6em,itemsep=7pt,topsep=6pt]
  \item $z+w$ \qquad and \qquad $z-w$ \qquad and \qquad $w-z$
  \item $zw$, twice: once by expanding the binomials and using $i^{2}=-1$, and
        once by substituting into the formula
        $(a+bi)(c+di)=(ac-bd)+(ad+bc)i$. The two must agree.
  \item $z^{2}$
  \item $\conj{z}$, and $z\conj{z}$
  \item $\dfrac{z}{w}$, by multiplying through by $\conj{w}/\conj{w}$
  \item $\dfrac{1}{w}$
  \item $\lvert z\rvert$, $\lvert w\rvert$, and $\lvert zw\rvert$.
        Compare $\lvert zw\rvert$ with $\lvert z\rvert\cdot\lvert w\rvert$ and
        state what you notice. (You will prove it in Problem D.3.)
\end{enumerate}
\end{prob}

\begin{prob}[Powers of $i$]
Compute, in the form $a+bi$ with $a$ and $b$ real:
\[
  i^{3},\qquad i^{4},\qquad i^{5},\qquad i^{10},\qquad i^{27},\qquad i^{100},
  \qquad (-i)^{3},\qquad (2i)^{4}.
\]
The first seven each come out to one of $1$, $-1$, $i$, $-i$. The last one does
not --- say why not, and what the extra factor is.

\noindent Then: explain in one or two sentences why the powers of $i$ repeat with
period $4$, and describe what that repetition looks like on the complex plane.
Finally, give a rule that computes $i^{n}$ for any positive integer $n$ from the
remainder of $n$ on division by $4$.
\end{prob}

\begin{prob}[The conjugate]
Let $z=5-12i$.

\begin{enumerate}[label=\textbf{(\alph*)},leftmargin=2.6em,itemsep=7pt,topsep=6pt]
  \item Compute $\conj{z}$, $z\conj{z}$, and $\lvert z\rvert$. Verify that
        $\lvert z\rvert^{2}=z\conj{z}$.
  \item Compute $z+\conj{z}$ and $z-\conj{z}$. Express each in terms of
        $\Rez(z)$ and $\Imz(z)$, and guess the general identity. (You will prove
        it in Problem D.1.)
  \item Plot $z$ and $\conj{z}$ and describe the geometric relationship between
        them.
  \item Compute $\conj{\conj{z}}$. What does conjugating twice do, and why is
        that obvious from the picture?
  \item Find every complex number $z$ satisfying $z=\conj{z}$. Describe the
        answer as a subset of the complex plane.
\end{enumerate}
\end{prob}

\begin{prob}[Modulus, distance, and discs]
\begin{enumerate}[label=\textbf{(\alph*)},leftmargin=2.6em,itemsep=7pt,topsep=6pt]
  \item Compute $\lvert -3+4i\rvert$, $\lvert 1+i\rvert$, $\lvert -7\rvert$,
        $\lvert 2i\rvert$, and $\lvert 0\rvert$.
  \item Compute the distance from $3+2i$ to $-1-i$, and the distance from
        $5-i$ to $5+3i$. For the second, explain why you could have read the
        answer off the picture without computing.
  \item Describe in words, and sketch, the set of complex numbers satisfying
        each condition:
        \[
          \lvert z\rvert < 2,
          \qquad
          \lvert z - (2-i)\rvert = 3,
          \qquad
          \lvert z - 4i\rvert \le 1 .
        \]
        Say for each whether the boundary circle is included.
  \item Write down an inequality describing the closed disc of radius $5$
        centred at $-3+2i$.
  \item Does $1+i$ lie in the disc $\lvert z-(2-i)\rvert\le 3$? Decide by
        computing, and say how far inside or outside it is.
  \item Describe the set $\lvert z\rvert = 0$. How many complex numbers are in
        it?
\end{enumerate}
\end{prob}

\begin{prob}[The unit circle, trigonometric form, and Euler]
\begin{enumerate}[label=\textbf{(\alph*)},leftmargin=2.6em,itemsep=7pt,topsep=6pt]
  \item By computing the modulus of each, verify that all four of the complex
        numbers
        \[
          i,
          \qquad
          \frac{\sqrt2}{2}+\frac{\sqrt2}{2}\,i,
          \qquad
          \frac12+\frac{\sqrt3}{2}\,i,
          \qquad
          -\frac35+\frac45\,i
        \]
        lie on the unit circle.
  \item Write each of the four in trigonometric form $\cos\theta+i\sin\theta$,
        with $\theta\in[0,2\pi)$. The first three have exact angles; for the
        fourth, give $\theta$ to four decimal places (a calculator is fine, or
        use \texttt{np.angle} from Part E) and say which quadrant it is in.
  \item Use Euler's formula to evaluate, in the form $a+bi$:
        \[
          e^{i\pi/2},\qquad
          e^{i\pi},\qquad
          e^{2\pi i},\qquad
          e^{i\pi/4},\qquad
          e^{-i\pi/2}.
        \]
  \item State Euler's identity, and say which of your answers in (c) it is a
        rearrangement of.
  \item A point of the unit circle needs only \emph{one} real number to pin it
        down, where a general complex number needs two. Explain that in one or
        two sentences, and say what the missing second number is being fixed to.
\end{enumerate}
\end{prob}

\begin{prob}[Multiplication by $i$ as a rotation]
Let $z = 3+2i$.

\begin{enumerate}[label=\textbf{(\alph*)},leftmargin=2.6em,itemsep=7pt,topsep=6pt]
  \item Compute $iz$, $i^{2}z$, $i^{3}z$ and $i^{4}z$, and plot all five points
        $z, iz, i^{2}z, i^{3}z, i^{4}z$ on one complex plane. Join them in order.
        What shape do the four distinct points form?
  \item Compute $\lvert iz\rvert$ and compare it with $\lvert z\rvert$. Why is
        this worth checking, rather than stopping once the angle is known to be
        $90^{\circ}$?
  \item Compute $(-i)z$. Which way did the point turn this time?
  \item A point sits at $(4,1)$ in the plane. Encode it as a complex number and
        rotate it $90^{\circ}$ counter-clockwise about the origin by a single
        multiplication. Give the coordinates of the rotated point.
  \item Rotate the same point by $180^{\circ}$, again by a single multiplication.
        Which complex number do you multiply by, and why?
  \item Take an arbitrary $z=a+bi$ and compute $iz$ symbolically. Write both $z$
        and $iz$ as vectors in $\R^{2}$. The full argument finishes with a dot
        product --- an operation we meet in Part 2 of the vector unit.
        \emph{Without} using it, describe in words the relationship between
        $[\,a,\,b\,]$ and $[\,-b,\,a\,]$: what happened to the two coordinates,
        and what does that do to the arrow?
\end{enumerate}
\end{prob}

%==============================================================================
\section{Complex numbers --- proof}
%==============================================================================
\emph{Complex Numbers.} In every problem, write $z=a+bi$ and $w=c+di$ with
$a,b,c,d\in\R$ and compute. As in Part B, name the fact about real numbers that
each step relies on when it is not obvious.

\begin{prob}[Recovering the parts from the conjugate]
Let $z\in\C$. Prove:
\begin{enumerate}[label=\textbf{(\alph*)},leftmargin=2.6em,itemsep=7pt,topsep=6pt]
  \item $z+\conj{z} = 2\Rez(z)$;
  \item $z-\conj{z} = 2i\,\Imz(z)$;
  \item $z$ is real if and only if $z=\conj{z}$;
  \item $\Rez(z)=0$ if and only if $z=-\conj{z}$.
\end{enumerate}

\noindent Parts (c) and (d) are ``if and only if'' statements, so each needs
\emph{both} directions argued. Say explicitly which direction you are proving as
you prove it.
\end{prob}

\begin{prob}[Conjugation respects the arithmetic]
Let $z,w\in\C$. Prove:
\begin{enumerate}[label=\textbf{(\alph*)},leftmargin=2.6em,itemsep=7pt,topsep=6pt]
  \item $\conj{z+w} = \conj{z}+\conj{w}$;
  \item $\conj{zw} = \conj{z}\,\conj{w}$.
\end{enumerate}

\noindent For (b), expand both sides using the product formula and compare real
and imaginary parts. Then state, in one sentence, what (b) says
geometrically: if reflecting in the real axis is what conjugation \emph{does},
what have you just shown about reflecting a product versus multiplying two
reflections?
\end{prob}

\begin{prob}[The modulus is multiplicative]
Prove that $\lvert zw\rvert = \lvert z\rvert\,\lvert w\rvert$ for all
$z,w\in\C$.

\begin{enumerate}[label=\textbf{(\alph*)},leftmargin=2.6em,itemsep=7pt,topsep=6pt]
  \item Show $\lvert zw\rvert^{2} = (ac-bd)^{2}+(ad+bc)^{2}$ and expand it fully.
  \item Show that the expansion factors as $(a^{2}+b^{2})(c^{2}+d^{2})$. Note
        which two terms cancel --- that cancellation is the entire content of the
        identity.
\end{enumerate}
\end{prob}

\begin{prob}[The unit circle really is the unit circle]
Let $\theta$ be any real number. Prove:
\begin{enumerate}[label=\textbf{(\alph*)},leftmargin=2.6em,itemsep=7pt,topsep=6pt]
  \item $\lvert\cos\theta + i\sin\theta\rvert = 1$. Name the trigonometric
        identity you use.
  \item $\lvert e^{i\theta}\rvert = 1$, deduced from (a) and Euler's formula.
\end{enumerate}
\end{prob}

\begin{prob}[Multiplying by a point of the unit circle preserves distance]
Let $z\in\C$ and let $\theta\in\R$.

\begin{enumerate}[label=\textbf{(\alph*)},leftmargin=2.6em,itemsep=7pt,topsep=6pt]
  \item Prove that $\lvert z\,e^{i\theta}\rvert = \lvert z\rvert$.
        \hint{Problem D.3 plus Problem D.4(b). This is a two-line proof.}
  \item Deduce as a special case that $\lvert iz\rvert = \lvert z\rvert$ for all
        $z$, by identifying the value of $\theta$ for which $e^{i\theta}=i$.
\end{enumerate}
\end{prob}

%==============================================================================
\section{NumPy}
%==============================================================================
\emph{Required.}

\medskip
\noindent This part is a tutorial, then a specification of the two files you hand
in, then the problems. Read the tutorial sections typing the code in as you go,
read \emph{How to structure your submission} before you write anything you intend
to submit, then do Problems E.1--E.8. If you have never programmed before, that
is expected and fine: everything you need is below, and nothing here is harder
than the arithmetic you did by hand in Parts A and C.

\medskip
\noindent Two references live alongside this assignment in the course
repository: \textbf{Setup Guide 2} (installing Python, pip, and an IDE) and the
\textbf{Python Recipe Book} (\texttt{lectures/python-recipe/}), which is a
lookup table of the calls covered here and in later lectures. Use the recipe
book as a dictionary, not a textbook.

\subsection{Getting set up}

Install once:

\begin{lstlisting}
python -m pip install numpy
\end{lstlisting}

\noindent Every block below assumes this line at the top of your file:

\begin{lstlisting}
import numpy as np
\end{lstlisting}

\checkpoint{Create \texttt{hw1.py}, put the import line in it, add
\texttt{print(np.array([1.0, 2.0]))}, and run it. You should see
\texttt{[1. 2.]}. If you get \texttt{ModuleNotFoundError: No module named
'numpy'}, the install went to a different Python than the one running your file
--- Setup Guide 2 \S8 covers this exact failure.}

\subsection{Vectors: making them, and reading them}

A vector in $\R^{2}$ is a NumPy array of two numbers.

\begin{lstlisting}
import numpy as np

u = np.array([3.0, -4.0])       # the vector [3, -4]
v = np.array([-1.0, 2.0])       # the vector [-1, 2]

print(u[0])                     # 3.0   -- the first component
print(u[1])                     # -4.0  -- the second
print(len(u))                   # 2     -- the n in R^n
print(u.shape)                  # (2,)
\end{lstlisting}

\heads{\textbf{Indices start at zero.} The first component of $\vc{u}$ is
\texttt{u[0]}, not \texttt{u[1]}. Lecture counts from $1$; Python counts from
$0$. Subtract one.}

\heads{\textbf{Write \texttt{3.0}, not \texttt{3}, when the answers will not be
whole numbers.} An array built from integers stays an integer array, and some
operations will surprise you. Using floats throughout this assignment costs
nothing and avoids the issue entirely.}

\subsection{Addition, subtraction, and scaling}

NumPy spells these exactly as the lecture does.

\begin{lstlisting}
u + v           # array([ 2., -2.])   -- componentwise addition
u - v           # array([ 4., -6.])   -- componentwise subtraction
-3 * u          # array([-9., 12.])   -- scalar multiplication
-u              # array([-3.,  4.])   -- the mirror image, (-1)u
2*u - 3*v       # array([ 9., -14.])  -- combinations, as in Problem A.5
\end{lstlisting}

\heads{\textbf{A Python list is not a vector.} Ordinary square brackets with no
\texttt{np.array} around them make a \emph{list}, and \texttt{+} on two lists
means \emph{glue them end to end}:}

\begin{lstlisting}
[1, 2] + [3, 4]                     # [1, 2, 3, 4]   -- concatenation!
np.array([1, 2]) + np.array([3, 4]) # array([4, 6])  -- what you wanted
\end{lstlisting}

\noindent{\small\sffamily Every vector in this assignment must be wrapped in
\texttt{np.array}.}

\heads{\textbf{\texttt{u * v} is not a vector operation we have defined.} NumPy
will happily compute it --- it multiplies entry by entry --- but that is not
addition, not scaling, and not any operation from the lectures. Do not use it in
this assignment.}

\subsection{The norm}

\begin{lstlisting}
np.linalg.norm(u)                   # 5.0  -- the Euclidean norm ||u||
np.sqrt(u[0]**2 + u[1]**2)          # 5.0  -- the same thing, by Pythagoras
np.linalg.norm(u - v)               # the distance between the two heads
\end{lstlisting}

\noindent The name is worth parsing: \texttt{np} is NumPy, \texttt{linalg} is
its linear-algebra module, \texttt{norm} is the function. Read it as ``NumPy's
linear algebra norm''.

\subsection{Comparing floating-point numbers}

This is the one piece of programming discipline this assignment insists on.

\begin{lstlisting}
0.1 + 0.2 == 0.3                    # False   (!)
0.1 + 0.2                           # 0.30000000000000004
np.isclose(0.1 + 0.2, 0.3)          # True    -- the right way to compare
\end{lstlisting}

\noindent Computers store real numbers to about sixteen significant digits, so
two quantities that are mathematically equal can differ in the last few. Testing
them with \texttt{==} then reports \texttt{False} and tells you nothing.

\begin{callout}{The rule for this course}
Never compare computed floating-point values with \texttt{==}. Use
\texttt{np.isclose(x, y)} for two numbers, and \texttt{np.allclose(a, b)} for two
arrays. Both work on complex numbers too.
\end{callout}

\subsection{Complex numbers in Python}

\begin{lstlisting}
z = 3 + 4j              # NOT 3 + 4i -- Python writes the imaginary unit as j
w = complex(1, -2)      # the same as 1 - 2j

z.real                  # 3.0            -- Re(z), a real number
z.imag                  # 4.0            -- Im(z), a real number
z.conjugate()           # (3-4j)         -- the conjugate, z-bar
abs(z)                  # 5.0            -- the modulus |z|

z + w                   # (4+2j)
z * w                   # (11-2j)
z / w                   # (-1+2j)
z ** 2                  # (-7+24j)

np.angle(z)             # 0.9272952...   -- the angle theta, in radians
np.exp(1j * np.pi / 2)  # ~ (6.1e-17+1j) -- Euler's formula: e^{i pi/2} = i
\end{lstlisting}

\heads{\textbf{The imaginary unit is \texttt{j}, and it must carry a
coefficient.} Writing \texttt{j} alone is a variable name and will raise
\texttt{NameError}; the imaginary unit is \texttt{1j}. (Python borrowed
\texttt{j} from electrical engineering, where $i$ already means current.) Also
note \texttt{1j**2} prints as \texttt{(-1+0j)}, not \texttt{-1} --- the value is
$-1$, but the type stays complex.}

\heads{\textbf{\texttt{z.imag} is a real number.} For \texttt{z = 3 + 4j} it is
\texttt{4.0}, not \texttt{4j} --- which is exactly the convention used in
lecture. \texttt{.real} and \texttt{.imag} are attributes, so they take no
parentheses; \texttt{.conjugate()} is a method, so it does.}

\heads{\textbf{The last line of the block above prints
\texttt{(6.123233995736766e-17+1j)}, not a clean \texttt{1j}.} That
$6.1\times10^{-17}$ is floating-point error, not mathematics: $\cos(\pi/2)$ is
zero but the stored value of $\pi$ is not exactly $\pi$. This is why
\texttt{np.isclose} exists, and Problem E.6 asks you to think about it.}

\subsection{How to structure your submission}

Part E is submitted as \textbf{two} Python files with different jobs.

\begin{itemize}[leftmargin=1.4em,itemsep=6pt,topsep=4pt]
  \item \texttt{hw1.py} --- \textbf{your work}. One function per lettered part,
        named \texttt{q}\,\emph{number}\,\emph{letter}: Problem E.2(a) becomes
        \texttt{q2a}, Problem E.6(c) becomes \texttt{q6c}. A question with no
        lettered parts gets one function named for the question alone, so E.8
        becomes \texttt{q8}. Each function \textbf{returns} its answer; when an
        answer has several pieces, return them together as a tuple or a list.
  \item \texttt{runhw1.py} --- \textbf{the driver}. It imports those functions
        from \texttt{hw1.py} and calls them in order, one line per part, printing
        a label before each answer. It contains no mathematics of its own.
        \textbf{This is the file that is run to grade Part E.}
\end{itemize}

\medskip
\noindent\textbf{\texttt{hw1.py} --- definitions only.} Nothing may run when the
file is imported: define your functions and stop. A \texttt{print} sitting at the
left margin fires the moment \texttt{runhw1.py} imports the file, and its output
lands above the labels, where it cannot be matched to a question.

\begin{lstlisting}
# hw1.py -- the work.  Definitions only; nothing runs on import.

import numpy as np

u = np.array([3.0, -4.0])
v = np.array([-1.0, 2.0])

def norm2(vec):                 # helpers are fine, and may be shared
    return np.sqrt(vec[0]**2 + vec[1]**2)

# --- E.2(a) -------------------------------------------------------
def q2a():
    return norm2(u)

# --- E.2(b) -------------------------------------------------------
def q2b():
    return [bool(np.isclose(norm2(w), np.linalg.norm(w))) for w in (u, v)]
\end{lstlisting}

\noindent\textbf{\texttt{runhw1.py} --- one line per part.} Import the functions
from their source file, then call them in the order the questions appear. Each
line prints a label followed by the answer that call produced:

\begin{lstlisting}
# runhw1.py -- runs every question in order.  No mathematics in this file.

import hw1

print("Output of question 2 part a : ", hw1.q2a())
print("Output of question 2 part b : ", hw1.q2b())
print("Output of question 2 part c : ", hw1.q2c())
\end{lstlisting}

\noindent (\texttt{from hw1 import q2a, q2b, q2c} and then \texttt{q2a()} is
equally acceptable if you prefer it.)

\medskip
\noindent\textbf{The label is required, and its wording is fixed:}

\begin{center}
\texttt{"Output of question <number> part <letter> : "}
\end{center}

\noindent So Problem E.6(c) is labelled \texttt{Output of question 6 part c :}\ .
For a question with no lettered parts, drop that half ---
\texttt{"Output of question 8 : "}.

\heads{Several problems below say ``print'' something. Return it instead ---
\texttt{runhw1.py} does all the printing, and that is where the value shows up
under its label. Nothing in \texttt{hw1.py} should print.}

\checkpoint{Run \texttt{runhw1.py} once before you push and read the output top
to bottom. Every part of every question should appear, in order, under its own
label, with no error and nothing unlabelled.}

\subsection{The problems}

\begin{prob}[Checking Part A]
Enter $\vc{u}=[3,-4]$, $\vc{v}=[-1,2]$, $\vc{w}=[5,0]$ as NumPy arrays and print
the value of every one of the nine sub-parts (a)--(i) of Problem A.2. Compare
each against what you computed by hand and state, in a comment, whether they
agree. If any disagree, find the error --- it is far more often in the hand
computation than in the code, and finding it is the point of the exercise.
\end{prob}

\begin{prob}[Writing the norm yourself]
\texttt{np.linalg.norm} is a black box. Open it.

\begin{enumerate}[label=\textbf{(\alph*)},leftmargin=2.6em,itemsep=7pt,topsep=6pt]
  \item Write a function \texttt{norm2(v)} that takes a two-component array and
        returns $\sqrt{x^{2}+y^{2}}$, using only \texttt{np.sqrt}, indexing, and
        arithmetic --- not \texttt{np.linalg.norm}.
  \item Test it against \texttt{np.linalg.norm} on all four vectors of
        Problem A.1, using \texttt{np.isclose}, and return the four verdicts.
  \item Use \texttt{norm2} to give the norms of the four vectors from
        Problem A.1 as decimals, and check them against the exact values you
        wrote by hand.
\end{enumerate}

\begin{lstlisting}
def norm2(v):
    # your two lines here
    ...

# --- E.2(b) -------------------------------------------------------
def q2b():
    return [bool(np.isclose(norm2(w), np.linalg.norm(w))) for w in vectors]
\end{lstlisting}
\end{prob}

\begin{prob}[Head minus tail]
Write a function \texttt{vector\_from(tail, head)} that takes two points, given
as NumPy arrays, and returns the vector they represent --- head minus tail.

Use it on all four rows of the table in Problem A.1, return the four results, and
then use \texttt{np.allclose} to confirm in code which pair are the same vector.
Return that answer as a sentence, for example \texttt{"rows (x) and (y) are the
same vector: True"}.
\end{prob}

\begin{prob}[Testing the triangle inequality at scale]
By hand you checked four pairs. Check ten thousand.

\begin{enumerate}[label=\textbf{(\alph*)},leftmargin=2.6em,itemsep=7pt,topsep=6pt]
  \item Using the generator below, draw $10{,}000$ random pairs
        $\vc{u},\vc{v}\in\R^{2}$ and count how many satisfy
        $\norm{\vc{u}+\vc{v}}\le\norm{\vc{u}}+\norm{\vc{v}}$. Print the count.
  \item In the same loop, count how many satisfy it with \emph{equality}, in the
        \texttt{np.isclose} sense. Print that count too, and explain the number
        you get in one sentence.
  \item Now construct, by hand in code, one pair for which equality \emph{does}
        hold, and verify it with \texttt{np.isclose}. Say what it is about your
        pair that forces equality.
  \item Answer in a comment: your run in (a) found no counterexample. Does that
        \emph{prove} the triangle inequality? Say why or why not, in one or two
        sentences.
\end{enumerate}

\begin{lstlisting}
rng = np.random.default_rng(231)    # a fixed seed, so your run is reproducible

for _ in range(10000):
    u = rng.normal(size=2)          # a random vector in R^2
    v = rng.normal(size=2)
    # your test here
\end{lstlisting}
\end{prob}

\begin{prob}[Implementing complex arithmetic from the formulas]
Python already multiplies and divides complex numbers. Ignore that, and
implement the multiplication and division rules yourself.

Write three functions that take the four \emph{real} numbers $a,b,c,d$ and
return the answer as a Python complex number, using only real arithmetic
internally and the formulas from the lecture:

\begin{enumerate}[label=\textbf{(\alph*)},leftmargin=2.6em,itemsep=7pt,topsep=6pt]
  \item \texttt{cmul(a, b, c, d)}, computing $(a+bi)(c+di)$ by the boxed formula
        $(ac-bd)+(ad+bc)i$;
  \item \texttt{cdiv(a, b, c, d)}, computing $(a+bi)/(c+di)$ by the conjugate
        trick, and raising an error if $c$ and $d$ are both zero;
  \item \texttt{cmod(a, b)}, computing $\lvert a+bi\rvert$ from Pythagoras.
\end{enumerate}

\noindent Then test all three against Python's built-in \texttt{*}, \texttt{/}
and \texttt{abs} on the pair $z=3-2i$, $w=-1+4i$ from Problem C.2, using
\texttt{np.isclose}, and return the verdicts. Finally, call \texttt{cdiv} with
$c=d=0$ inside a \texttt{try} block, catch the error, and return its message, to
show the guard works.

\begin{lstlisting}
def cmul(a, b, c, d):
    return complex(..., ...)        # (ac - bd) + (ad + bc)i
\end{lstlisting}

\heads{Build the answer with \texttt{complex(real\_part, imag\_part)}. If you
find yourself writing \texttt{(a + b*1j) * (c + d*1j)} you have used Python's
multiplication instead of implementing the lecture's formula, which is the one
thing this problem is asking you not to do.}
\end{prob}

\begin{prob}[Euler's formula, numerically]
\begin{enumerate}[label=\textbf{(\alph*)},leftmargin=2.6em,itemsep=7pt,topsep=6pt]
  \item For each of the ten values $\theta = 0$, $\pi/10$, $2\pi/10$, \dots,
        $9\pi/10$, compute \texttt{np.exp(1j*theta)} and
        \texttt{np.cos(theta) + 1j*np.sin(theta)}, and confirm with
        \texttt{np.isclose} that they agree. Return the ten $\theta$ values
        paired with their ten verdicts. The block at the end of this problem
        builds all ten values of $\theta$ for you in one line.
  \item Print \texttt{np.exp(1j*np.pi) + 1}. It is not exactly zero. Print
        \texttt{np.isclose(np.exp(1j*np.pi) + 1, 0)} as well. In a comment,
        explain in two sentences why the exact answer is $0$ (Euler's identity)
        but the computed answer is not, and what the size of the discrepancy
        tells you.
  \item Confirm that $\lvert e^{i\theta}\rvert = 1$ for all ten values of
        $\theta$ from (a), using \texttt{np.abs} and \texttt{np.allclose}. This
        is Problem D.4(b), checked numerically --- note again that checking it on
        ten values is not what proving it means.
  \item Verify your answers to Problem C.6(b): for each of the four numbers
        there, print \texttt{np.angle} of it and check it against the $\theta$
        you found by hand.
\end{enumerate}

\begin{lstlisting}
thetas = np.pi * np.arange(10) / 10     # 0, pi/10, 2pi/10, ..., 9pi/10
\end{lstlisting}
\end{prob}

\begin{prob}[Rotation by $i$]
\begin{enumerate}[label=\textbf{(\alph*)},leftmargin=2.6em,itemsep=7pt,topsep=6pt]
  \item For \texttt{z = 3 + 2j}, print $i^{k}z$ for $k=0,1,2,3,4$ and check that
        the fifth value equals the first, using \texttt{np.isclose}. Compare with
        your hand answers to Problem C.7(a).
  \item Draw $1000$ random complex numbers and verify with \texttt{np.allclose}
        that $\lvert iz\rvert = \lvert z\rvert$ for every one of them.
        \hint{\texttt{zs = rng.normal(size=1000) + 1j*rng.normal(size=1000)}
        makes all thousand at once, and \texttt{np.abs(zs)} takes all thousand
        moduli at once --- no loop needed.}
  \item Take the four corners of the unit square, $0$, $1$, $1+i$, $i$, as an
        array. Multiply the whole array by $i$ and print the result. Which corner
        went where? Multiply by $i$ three more times and confirm you are back
        where you started.
  \item Rotate the point $(4,1)$ by $90^{\circ}$ counter-clockwise, by
        $180^{\circ}$, and by $\theta = \pi/6$, each with a single complex
        multiplication. Print the resulting coordinates as ordinary pairs of real
        numbers, using \texttt{.real} and \texttt{.imag}. Check the first two
        against your hand answers to Problem C.7(d) and (e).
\end{enumerate}
\end{prob}

\begin{prob}[Checking Part C]
Print the value of every sub-part of Problem C.2, of Problem C.4(a)--(b), and of
Problem C.5(a)--(b), and state in a comment whether each agrees with your hand
computation.

\medskip
\noindent Then answer, in a comment at the end of your file, in three or four
sentences: which problems on this assignment could NumPy check for you, and
which could it not? Name at least one problem from Part B or Part D and say
precisely what a computer run would and would not establish about it.
\end{prob}

\begin{callout}{What NumPy can and cannot do for you}
\renewcommand{\arraystretch}{1.25}
\begin{tabular}{@{}l l@{}}
\toprule
\textbf{Part} & \textbf{Can NumPy check it?} \\
\midrule
A --- vectors, computation      & Yes, entirely. Do it before you submit. \\
B --- vectors, proof            & No. It can test instances, which is not a proof. \\
C --- complex, computation      & Yes, entirely. Do it before you submit. \\
D --- complex, proof            & No, for the same reason. \\
E --- NumPy                     & It \emph{is} the check. \\
\bottomrule
\end{tabular}

\smallskip
For the drawing sub-parts of Problems A.1, C.1, C.4 and C.7, nothing is
being checked by machine --- but sketching the picture is often the fastest way
to notice that a sign is wrong.
\end{callout}

%==============================================================================
\section*{Optional and ungraded: drawing the pictures}
%==============================================================================
\noindent Not required, not graded, and not needed for any later assignment ---
but if you would like the computer to draw the figures you have been sketching
by hand, this is all it takes. It needs one more package:

\begin{lstlisting}
python -m pip install matplotlib
\end{lstlisting}

\heads{Keep this out of \texttt{hw1.py} and \texttt{runhw1.py}, in a scratch file
of its own. \texttt{plt.show()} halts a program until you close the window, so a
plot left in your submission stalls the run it is graded from.}

\noindent Two vectors and their sum, drawn as arrows from the origin:

\begin{lstlisting}
import numpy as np
import matplotlib.pyplot as plt

u = np.array([2.0, 1.0])
v = np.array([1.0, 3.0])

for vec, colour, name in [(u, "C0", "u"), (v, "C2", "v"), (u+v, "C3", "u+v")]:
    plt.arrow(0, 0, vec[0], vec[1], head_width=0.12,
              length_includes_head=True, color=colour)
    plt.text(vec[0], vec[1], f" {name}", color=colour)

plt.axhline(0, color="gray", lw=0.5)
plt.axvline(0, color="gray", lw=0.5)
plt.gca().set_aspect("equal")
plt.grid(alpha=0.3)
plt.show()
\end{lstlisting}

\noindent The four powers of $i$ times $z$, on the complex plane:

\begin{lstlisting}
z = 3 + 2j
points = np.array([z * 1j**k for k in range(4)])

plt.scatter(points.real, points.imag)
for p in points:
    plt.text(p.real, p.imag, f"  {p:.0f}")

plt.gca().set_aspect("equal")
plt.axhline(0, color="gray", lw=0.5)
plt.axvline(0, color="gray", lw=0.5)
plt.grid(alpha=0.3)
plt.show()
\end{lstlisting}

\noindent \texttt{points.real} and \texttt{points.imag} take the real and
imaginary parts of the whole array at once --- the same attributes as for a
single complex number, applied elementwise.

\bigskip
\noindent\rule{\textwidth}{0.4pt}

\smallskip
\noindent{\small Questions about any problem here belong in the class Discord,
where everyone can see the answer. Questions that are specific to your own work
--- ``is my proof of B.4 valid?'' --- are better brought to office hours,
Thursdays 3:00--4:00 PM in Kiely Hall 508. Please email ahead.}

\end{document}
